\section{Implementation details}

\subsection{Actor inputs}
The upper operational block uses phase sine and
cosine, remaining day fraction, each directional remaining-trip fraction
relative to the remaining day, quota-available gaps, directional ready
counts, total active fleet fraction, the cycle-based spacing floor, recent
budget utilization and preceding interval readiness delay. For each
direction, the four demand statistics are the sum, scaled spatial root
mean square, scaled spatial standard deviation and aggregate one-bin
change. For the lower actor the operational inputs are the dynamic action
limit, remaining holding, holding per remaining stop, estimated turnaround
slack, overdue trips at the next terminal, two relative headways, two-sided
availability, analytical gap-based hold target, the two planned target
headways, onboard fraction, route progress, and phase sine and cosine.
Its demand statistics are the local quantity and the mean of the next
two stops, each with its one-bin change. This summary and its normalization
are shared by the corrected input comparisons.

\subsection{Updates}
Each policy and critic has two 96-unit hidden layers.
The reward critics use six members and a pessimistic mean-minus-standard
deviation target. The upper replay uses six-step discounted returns
($\gamma=0.996$), a capacity of 20,000 and batches of 48. The lower replay
uses one-step returns ($\gamma=0.92$), a capacity of 180,000 and batches of
128. Both learning rates are $3\times10^{-4}$; the lower multiplier uses
$10^{-3}$ and a cost target of 0.020. The four final seed pairs use the
first two days as replay warm-up and make 32 upper and 64 lower updates
on each subsequent training day. Final evaluation uses greedy action
selection with the same action mask as training.

\subsection{Limits and forecasts}
The dynamic lower mask multiplies the
normalized remaining action allowance by the 25-second bin maximum;
the zero bin is always legal. The upper spacing floor estimates the
two directional trip times plus 40 seconds divided by the fleet size,
then scales by 0.85, 0.90 or 0.97 as terminal stock falls from at least
three to two or fewer than two vehicles. The result is clipped to
$[250,335]$ seconds, with an upper requested gap of 455 seconds. The
operating projection falls back to uniform residual gaps when the
two-block allocation is infeasible. The minimum actual dispatch
spacing enforced by the simulator is a separate 120-second rule.
